% =====================================================================
\chapter{Probabilidad, condicionamiento y teorema de Bayes}
\label{tema:1}
% =====================================================================

Este tema fija el lenguaje del módulo: espacios de probabilidad, probabilidad condicionada, independencia y teorema de Bayes. La presentación sigue el capítulo~1 de~\cite{bt} y las lecciones~1 a~3 de~\cite{mit6041}; la terminología bayesiana (a priori, verosimilitud, a posteriori) y la forma de razón de probabilidades siguen las clases~3 y~12 de~\cite{mit1805} y el capítulo~2 de~\cite{bh}. Deisenroth \emph{et al.}~\cite[§6.1 y §6.3]{mml} ofrecen el mismo material con orientación al aprendizaje automático.

\section{Modelos probabilísticos}

\begin{definicion}[Espacio de probabilidad]\label{def:espacio}
Un \emph{espacio de probabilidad} es una terna $(\Omega,\mathcal{F},\PP)$ formada por un \emph{espacio muestral} $\Omega$ (el conjunto de resultados posibles de un experimento), una familia $\mathcal{F}$ de subconjuntos de $\Omega$, llamados \emph{sucesos}, cerrada por complementos y uniones numerables y que contiene a $\Omega$, y una función $\PP\colon\mathcal{F}\to\R$ que cumple los \emph{axiomas de Kolmogorov}:
\begin{enumerate}
  \item \emph{no negatividad}: $\PP(A)\ge 0$ para todo $A\in\mathcal{F}$;
  \item \emph{normalización}: $\PP(\Omega)=1$;
  \item \emph{aditividad numerable}: si $A_1,A_2,\dots$ son disjuntos dos a dos, $\PP\bigl(\bigcup_{i}A_i\bigr)=\sum_{i}\PP(A_i)$.
\end{enumerate}
\end{definicion}

Cuando $\Omega$ es finito o numerable se toma $\mathcal{F}$ igual a todos los subconjuntos de $\Omega$, y basta dar la probabilidad de cada resultado: $\PP(A)=\sum_{\omega\in A}\PP(\{\omega\})$. Si además $\Omega$ es finito y todos los resultados son equiprobables, $\PP(A)=\abs{A}/\abs{\Omega}$ (ley de Laplace)~\cite[§1.2]{bt}.

\begin{proposicion}[Consecuencias de los axiomas]\label{prop:axiomas}
Para sucesos $A,B\in\mathcal{F}$:
\begin{enumerate}
  \item $\PP(\varnothing)=0$ y $\PP(A^{c})=1-\PP(A)$;
  \item si $A\subseteq B$, entonces $\PP(A)\le\PP(B)$;
  \item $\PP(A\cup B)=\PP(A)+\PP(B)-\PP(A\cap B)$;
  \item $\PP(A\cup B)\le\PP(A)+\PP(B)$ (cota de la unión).
\end{enumerate}
\end{proposicion}

\section{Probabilidad condicionada}

\begin{definicion}[Probabilidad condicionada]
Sean $A,B\in\mathcal{F}$ con $\PP(A)>0$. La \emph{probabilidad de $B$ condicionada a $A$} es
\[
\PP(B\mid A)=\frac{\PP(A\cap B)}{\PP(A)}.
\]
\end{definicion}

Condicionar a $A$ equivale a tomar $A$ como nuevo espacio muestral y renormalizar. Para $A$ fijo, $B\mapsto\PP(B\mid A)$ cumple los tres axiomas de Kolmogorov~\cite[§1.3]{bt}; en consecuencia, todas las propiedades de la \cref{prop:axiomas} valen también para probabilidades condicionadas. Despejando se obtiene la \emph{regla del producto}, $\PP(A\cap B)=\PP(A)\,\PP(B\mid A)$, cuya iteración es la regla de la cadena del módulo~1.

\begin{observacion}
Un modelo de lenguaje es una familia de probabilidades condicionadas $\PP(x_t\mid x_{1:t-1})$: el suceso condicionante es el contexto observado y la distribución condicionada es la del token siguiente. Todo lo que se demuestre en este tema para probabilidades condicionadas se aplica, por tanto, a las salidas de un LLM.
\end{observacion}

\section{Probabilidad total y regla de Bayes}

\begin{definicion}[Partición]
Una familia finita o numerable de sucesos $H_1,H_2,\dots$ es una \emph{partición} de $\Omega$ si son disjuntos dos a dos y su unión es $\Omega$.
\end{definicion}

\begin{teorema}[Probabilidad total]\label{teo:total}
Sea $H_1,H_2,\dots$ una partición de $\Omega$ con $\PP(H_i)>0$ para todo $i$. Para todo suceso $D$,
\[
\PP(D)=\sum_{i}\PP(D\mid H_i)\,\PP(H_i).
\]
\end{teorema}

\begin{teorema}[Regla de Bayes]\label{teo:bayes}
Con las hipótesis del \cref{teo:total} y $\PP(D)>0$, para cada $j$,
\[
\PP(H_j\mid D)=\frac{\PP(D\mid H_j)\,\PP(H_j)}{\sum_{i}\PP(D\mid H_i)\,\PP(H_i)}.
\]
\end{teorema}

En el lenguaje de la inferencia, los $H_i$ son \emph{hipótesis} y $D$ son los \emph{datos}. $\PP(H_j)$ es la probabilidad \emph{a priori}, $\PP(D\mid H_j)$ la \emph{verosimilitud}, $\PP(H_j\mid D)$ la probabilidad \emph{a posteriori} y el denominador, $\PP(D)$, la \emph{evidencia}, que solo normaliza. Abreviadamente, la probabilidad a posteriori es proporcional al producto de verosimilitud y probabilidad a priori~\cite[clase~11]{mit1805}.

\addtocounter{definicion}{1}%

\addtocounter{figure}{1}%

Cuando solo hay dos hipótesis, $H$ y $H^{c}$, es cómodo trabajar con la \emph{razón de probabilidades} (en inglés, \emph{odds}) $O(H)=\PP(H)/\PP(H^{c})$.

\begin{proposicion}[Bayes en forma de razón]\label{prop:odds}
Si $\PP(D\mid H^{c})>0$, entonces
\[
O(H\mid D)=\frac{\PP(D\mid H)}{\PP(D\mid H^{c})}\cdot O(H),
\]
donde $O(H\mid D)=\PP(H\mid D)/\PP(H^{c}\mid D)$. El cociente $\PP(D\mid H)/\PP(D\mid H^{c})$ se llama \emph{factor de Bayes} o razón de verosimilitudes.
\end{proposicion}

La forma de razón muestra con claridad la llamada \emph{falacia de la probabilidad a priori}: un factor de Bayes grande no basta para que la hipótesis sea probable si $O(H)$ es muy pequeña (\cref{eje:1-detector,eje:1-cribado}).

\section{Independencia}

\begin{definicion}[Independencia]\label{def:indep}
Dos sucesos $A$ y $B$ son \emph{independientes} si $\PP(A\cap B)=\PP(A)\,\PP(B)$. Si $\PP(A)>0$, la condición equivale a $\PP(B\mid A)=\PP(B)$: saber que ha ocurrido $A$ no cambia la probabilidad de $B$. Los sucesos $A_1,\dots,A_n$ son \emph{independientes} si $\PP\bigl(\bigcap_{i\in S}A_i\bigr)=\prod_{i\in S}\PP(A_i)$ para todo subconjunto $S\subseteq\{1,\dots,n\}$.
\end{definicion}

\begin{definicion}[Independencia condicionada]
Sea $C$ con $\PP(C)>0$. Los sucesos $A$ y $B$ son \emph{independientes condicionalmente a $C$} si $\PP(A\cap B\mid C)=\PP(A\mid C)\,\PP(B\mid C)$.
\end{definicion}

\begin{observacion}
La independencia condicionada no implica la independencia, ni al revés. Dos lanzamientos de una moneda elegida al azar entre una equilibrada y una trucada son independientes condicionalmente a la moneda, pero no lo son incondicionalmente: una cara en el primer lanzamiento hace más probable la moneda que favorece la cara y, con ella, una cara en el segundo (\cref{eje:1-monedas}). Del mismo modo, independencia no es lo mismo que incompatibilidad: dos sucesos disjuntos con probabilidad positiva nunca son independientes (\cref{eje:1-disjuntos}).
\end{observacion}

\addtocounter{definicion}{1}%

